% Image credits for the photographs and AI illustrations of Book 2 (High School), Indonesian edition.
% Machine-readable license records live in images/book2/CREDITS.md; keep this
% file in sync with frontmatter/image-credits-book2.tex (one \item per item there).
% Names, institutions, datasets, image identifiers and licences are kept exactly as in
% the English file; only the frame and the descriptions are translated.
\chapter*{Kredit Gambar}

Gambar, peta, dan grafik dalam buku ini adalah karya asli; peta dan grafik
yang digambar dari data yang telah diterbitkan menyebutkan sumbernya di
keterangan gambar atau di daftar ini. Foto-foto dipakai di bawah lisensi
bebasnya masing-masing, dengan terima kasih kepada para pembuatnya; tautan
sumber lengkap dan alamat lisensi setiap foto tercantum dalam
\texttt{images/book2/CREDITS.md} di repositori buku ini.

\bigskip
Selain foto dan gambar asli, sebagian gambar adalah ilustrasi buatan AI,
yang dihasilkan dengan pembangkit gambar OpenAI dari perintah yang ditulis
oleh para penyunting buku ini, dan diperiksa ketepatan ilmiahnya sebelum
dimuat. Ilustrasi itu hanya menampilkan adegan umum, tidak pernah tempat
yang nyata, orang yang nyata, atau spesimen yang nyata. Perintah lengkap
untuk setiap gambar buatan AI dicatat dalam
\texttt{images/book2/ai/PROMPTS.md} di repositori.

\bigskip
\noindent Foto dan data, menurut bab:
\begin{itemize}\small
  \item Bab~1 (mengukur Bumi): Bumi dilihat dari Apollo~17; foto NASA, domain publik (Wikimedia Commons).
  \item Bab~1: sebuah lingkaran pengulang, pelat dari J.-B. Delambre, \emph{Base du système métrique décimal}, vol.~2 (1807); domain publik (Wikimedia Commons, dari BnF Gallica).
  \item Bab~1: panjang derajat dihitung pada elipsoid WGS~84 (set data EPSG).
  \item Bab~2 (planet yang layak huni): Venus dari Mariner~10 (PIA23791), NASA/JPL-Caltech, domain publik; Gunung Sharp dari wahana penjelajah Curiosity (PIA16769), NASA/JPL-Caltech/MSSS, domain publik (Wikimedia Commons).
  \item Bab~2: data planet dari lembar fakta NASA (NASA fact sheets); batas fase air dari terbitan IAPWS SR1-86 dan R14-08; batas zona layak huni dari Kopparapu dan kawan-kawan (2013); Proxima~b dari NASA Exoplanet Archive.
  \item Bab~3 (neraca energi Bumi): kolam lelehan di atas es laut Arktik (misi ICESCAPE, 2011); foto oleh NASA Goddard Space Flight Center, domain publik (Wikimedia Commons).
  \item Bab~3: nilai neraca energi dari Laporan Penilaian Keenam IPCC, Kelompok Kerja~I, Gambar~7.2; albedo dari University of Colorado Boulder (MOSAiC) dan National Snow and Ice Data Center; transpor panas dari Trenberth dan Caron (2001); produksi primer bersih dari Field dan kawan-kawan (1998); data termokimia dari NIST Chemistry WebBook.
  \item Bab~4 (mengapa awan terbentuk): kumulus cuaca cerah, foto karya Ralph F. Kresge, NOAA; landasan sebuah kumulonimbus, NOAA National Severe Storms Laboratory; keduanya domain publik (Wikimedia Commons).
  \item Bab~4: kejenuhan dihitung dari terbitan IAPWS; ukuran tetes dari NASA Jet Propulsion Laboratory.
  \item Bab~5 (badai dan siklon tropis): Badai Isabel dari Stasiun Antariksa Internasional, foto karya Mike Trenchard, NASA Johnson Space Center; sebuah tornado di Kansas, foto karya Sean Waugh, NOAA National Severe Storms Laboratory; petir, foto karya C.~Clark, NOAA National Severe Storms Laboratory; semuanya domain publik (Wikimedia Commons).
  \item Bab~5: lintasan siklon dari IBTrACS versi~4r01 (Gahtan dan kawan-kawan, 2024; Knapp dan kawan-kawan, 2010; NOAA National Centers for Environmental Information); kategori badai dari NOAA National Hurricane Center; energi badai dan syarat pembentukannya dari NOAA Atlantic Oceanographic and Meteorological Laboratory; laju kilat dari Christian dan kawan-kawan (2003); fakta tornado dari NOAA Storm Prediction Center.
  \item Bab~6 (samudra yang berlapis): sebuah roset pengambil sampel, foto karya Pablo Clemente-Colon, NOAA Photo Library; sebuah pelampung pemrofil yang diluncurkan dari kapal riset, NOAA Global Ocean Monitoring and Observing Program; klorofil samudra 2002--2010, Jesse Allen dan Robert Simmon, NASA Earth Observatory; semuanya domain publik (Wikimedia Commons).
  \item Bab~6: profil samudra dari World Ocean Atlas 2023 (Reagan dan kawan-kawan, 2024; NOAA National Centers for Environmental Information); massa jenis air laut dari persamaan keadaan internasional (Fofonoff dan Millard, 1983); volume samudra dari Charette dan Smith (2010); fakta pelampung dari Argo Program Office; upwelling dari NOAA National Ocean Service dan Kessler (2006).
  \item Bab~7 (tanah, sumber daya yang lambat terbentuk): profil laterit di atas gneis di Kerala, foto karya Werner Schellmann, CC~BY-SA~2.5; saprolit di Madagaskar tengah, foto karya Diorit, CC~BY-SA~3.0; sebuah profil tanah, USDA Natural Resources Conservation Service, domain publik (semuanya Wikimedia Commons).
  \item Bab~7: laju produksi dan erosi tanah dari Montgomery (2007); rumus mineral dari Mineralogy Database (webmineral.com); tebal saprolit dari USGS Ground Water Atlas of the United States.
  \item Bab~8 (daerah aliran sungai, sungai, dan banjir): seorang hidrolog mengukur debit sungai dan sebuah stasiun pengukur, U.S. Geological Survey, domain publik (Wikimedia Commons).
  \item Bab~8: puncak tahunan dan debit harian Sungai Potomac dekat Washington, U.S. Geological Survey National Water Information System, stasiun 01646500; debit dan daerah aliran Mississippi dari Kammerer (1990, USGS); debit Amazon dari Callède dan kawan-kawan (2010); keselamatan di air banjir dari NOAA National Weather Service.
  \item Bab~9 (sumber daya air tawar): Laut Aral pada 2000 dan 2018, NASA Earth Observatory (MODIS/Terra); ladang beririgasi dekat Sublette, Kansas, NASA (ASTER/Terra), domain publik (Wikimedia Commons).
  \item Bab~9: cadangan air dunia menurut Shiklomanov (1993), U.S. Geological Survey Water Science School; siklus air global menurut Trenberth dan kawan-kawan (2007), sebagaimana dicetak dalam Trenberth dan Asrar (2014); umur air tanah Nubia dari Sturchio dan kawan-kawan (2004); perubahan akuifer High Plains dari McGuire (2017, USGS); neraca Laut Aral dari Gaybullaev dan kawan-kawan (2012) dan Micklin (2007).
  \item Bab~10 (dari pergeseran benua sampai pemekaran dasar samudra): kapal pengebor \emph{JOIDES Resolution}, Ocean Drilling Program, Texas A\&M University, domain publik (Wikimedia Commons).
  \item Bab~10: skala waktu polaritas Cande dan Kent (1995); hasil pengeboran Leg~3 Deep Sea Drilling Project (1970); umur dasar samudra dari Seton dan kawan-kawan (2020), EarthByte, lisensi CC~BY~4.0.
  \item Bab~11 (gelombang seismik menyingkap bagian dalam Bumi): pemasangan seismometer sementara di Puerto Riko; foto karya Thomas Pratt, U.S. Geological Survey, domain publik (Wikimedia Commons).
  \item Bab~11: Preliminary Reference Earth Model dari Dziewonski dan Anderson (1981); waktu tempuh empiris gelombang P dan S dari Australian Passive Seismic Server (AuSPASS).
  \item Bab~12 (litosfer dan astenosfer): sebuah gunung es, NOAA, domain publik (USGS Water Science School).
  \item Bab~12: aliran panas dari Davies dan Davies (2010), lisensi CC~BY~3.0; provinsi aliran panas dari Roy, Blackwell dan Decker (1972); solidus mantel dari Hirschmann (2000); tebal dan massa jenis kerak dari Christensen dan Mooney (1995); ketinggian dari ETOPO1 (NOAA).
  \item Bab~13 (lempeng dan geraknya): sebuah stasiun penentuan posisi satelit di Pegunungan Cascade; foto karya Dave Sherrod, U.S. Geological Survey, domain publik (Wikimedia Commons). Relief berbayang dasar Pasifik Utara, National Geophysical Data Center, NOAA, domain publik (Wikimedia Commons).
  \item Bab~13: kutub rotasi lempeng NNR-MORVEL56 dari Argus, Gordon dan DeMets (2011); posisi stasiun dari Nevada Geodetic Laboratory (Blewitt, Hammond dan Kreemer, 2018); umur gunung api dari Clague dan Dalrymple (1987), U.S. Geological Survey; batas lempeng dari Bird (2003); daratan dari Natural Earth.
  \item Bab~14 (punggung samudra): lava bantal di Celah Galápagos, NOAA Okeanos Explorer Program, domain publik; sebuah cerobong hitam di Punggung Juan de Fuca, University of Washington dan NOAA, domain publik; gabro, foto karya Darla Sondrol (GeoDIL), CC0; peridotit di dalam basal, foto karya James St.~John, CC~BY~2.0 (semuanya Wikimedia Commons).
  \item Bab~14: kedalaman dasar samudra dari ETOPO1 (NOAA) dan umurnya dari Seton dan kawan-kawan (2020), EarthByte, lisensi CC~BY~4.0; hukum kedalaman dari Parsons dan Sclater (1977); solidus mantel dari Hirschmann (2000); kutub rotasi NNR-MORVEL56 dari Argus, Gordon dan DeMets (2011).
  \item Bab~15 (zona subduksi): sekis biru dan eklogit, foto karya James St.~John, CC~BY~2.0; kerusakan akibat tsunami di Kesennuma, Jepang, 2011, foto karya Dylan McCord, U.S. Navy, domain publik (semuanya Wikimedia Commons).
  \item Bab~15: gempa bumi dari katalog gempa USGS (ComCat); batas lempeng dan laju konvergensi dari Bird (2003); umur dasar samudra dari Seton dan kawan-kawan (2020), EarthByte, lisensi CC~BY~4.0; kecepatan lempeng NNR-MORVEL56 dari Argus, Gordon dan DeMets (2011); daratan dari Natural Earth.
  \item Bab~16 (mineral): kuarsa, ortoklas, biotit, dan olivin, foto karya Shannon Heinle, Darla Sondrol, Dexter Perkins dan Nessa Eull (Geology Digital Image Library), CC0; basal dalam sayatan tipis, foto karya Chmee2, CC~BY-SA~3.0; kuarsa dalam granit dalam sayatan tipis, foto karya Michael C.~Rygel, CC~BY-SA~3.0 (semuanya Wikimedia Commons).
  \item Bab~16: komposisi kerak benua dari Rudnick dan Gao (2014); massa atom standar IUPAC/CIAAW; rumus dan massa jenis mineral dari webmineral.com; massa jenis batuan granitik Sierra Nevada dari U.S. Geological Survey (Open-File Report 91-570-A).
  \item Bab~17 (batuan beku): kolom-kolom basal Giant's Causeway, foto karya Adam Bishop, CC~BY-SA~3.0; El Capitan, Yosemite, U.S. Geological Survey, CC0; basal, andesit, riolit, diorit, granit, obsidian, sebuah batuan porfiritik, dan dua sayatan tipis, foto karya Nessa Eull, Twyla Baker, Joseph H.~Hartman, Mike Davis, Dexter Perkins dan Darla Sondrol (Geology Digital Image Library), CC0 (semuanya Wikimedia Commons).
  \item Bab~17: pendinginan retas dan sill menurut Delaney (1986), U.S. Geological Survey Open-File Report 86-444; komposisi magma dan urutan kristalisasi dari catatan geologi fisik S.~A.~Nelson dan S.~Earle, \emph{Physical Geology} (CC~BY~4.0); biografi Bowen dari sejarah Geophysical Laboratory yang disusun Carnegie Science.
  \item Bab~18 (batu bara, minyak bumi, dan gas): fosil tumbuhan dalam serpih kaya karbon, lignit, batu bara bitumen, dan antrasit, foto karya James St.~John, CC~BY~2.0; sebuah lapisan lignit di Rumania, U.S. Geological Survey, domain publik (semuanya Wikimedia Commons).
  \item Bab~18: jendela minyak, pembentukan batu bara, dan perangkap dari S.~Earle, \emph{Physical Geology} (CC~BY~4.0); kandungan karbon batu bara dari U.S. Energy Information Administration; sistem minyak bumi menurut Magoon dan Schmoker (U.S. Geological Survey, 2000); cekungan rekahan dan cekungan garam di Afrika Tengah bagian barat dari U.S. Geological Survey Fact Sheet 2023-3028.
  \item Bab~19 (kerak benua): gneis Acasta, foto karya Pedroalexandrade, CC~BY-SA~3.0; zirkon Jack Hills, gambar karya A.~Cavosie dan J.~Valley, domain publik; gneis dan granit Perisai Kanada, foto karya Chris Light, CC~BY-SA~4.0 (semuanya Wikimedia Commons).
  \item Bab~19: komposisi kerak dari Rudnick dan Gao (2014); tebal kerak dari Christensen dan Mooney (1995), White, McKenzie dan O'Nions (1992) serta Karplus dan kawan-kawan (2019); luas benua dari Cogley (1984); umur dasar samudra dari Seton dan kawan-kawan (2020), EarthByte, lisensi CC~BY~4.0.
  \item Bab~20 (pembentukan pegunungan): Gunung Everest dari pesawat ulang-alik Columbia, NASA, domain publik; lapisan batuan terlipat di ngarai Kali Gandaki, foto karya Jean-Marie Hullot, CC~BY~2.0; Sesar Naik Glarus, foto karya Christian Heine, CC~BY-SA~3.0; lava bantal ofiolit Semail, foto karya Lionel Duchoiselle, CC~BY-SA~3.0 (semuanya Wikimedia Commons).
  \item Bab~20: posisi masa lalu India dan Lhasa dari model lempeng Müller dan kawan-kawan (2019, 2022), dihitung dengan GPlates Web Service (EarthByte); dataran tinggi dari ETOPO1; gempa bumi dari katalog USGS; batas lempeng dari Bird (2003); daratan dari Natural Earth.
  \item Bab~21 (batuan metamorf): belahan batu sabak di bawah mikroskop, foto karya Geoloshe, CC0; sekis mika bergarnet dari Syros, foto karya Graeme Churchard, CC~BY~2.0; sekis kianit dan gneis kuarsa--biotit, foto karya James St.~John, CC~BY~2.0 (semuanya Wikimedia Commons).
  \item Bab~21: medan kestabilan kianit, andalusit, dan silimanit dari Holdaway dan Mukhopadhyay (1993).
  \item Bab~22 (magmatisme di zona subduksi): kepulan letusan Gunung St.~Helens, 18 Mei 1980, foto karya Donald A. Swanson, USGS, domain publik; granodiorit, andesit, dan Sierra Nevada, foto karya James St.~John, CC~BY~2.0 (semuanya Wikimedia Commons).
  \item Bab~22: gempa bumi dari katalog USGS; gunung api dari Smithsonian Global Volcanism Program; batas lempeng dari Bird (2003); daratan dari Natural Earth; solidus peridotit kering dari Hirschmann (2000), dengan massa jenis PREM.
  \item Bab~23 (hancurnya pegunungan): Pegunungan Appalachia dari Stasiun Antariksa Internasional, foto astronaut ISS033-E-22378, NASA, domain publik; konglomerat kerakal kuarsa, foto karya James St.~John, CC~BY~2.0 (keduanya Wikimedia Commons).
  \item Bab~23: laju denudasi dari Portenga dan Bierman (2011); muatan sedimen Mississippi dari Meade dan Moody (2010); konstanta peluruhan dari Ahnert (1970); Perbukitan Varden, wilayah rekaan seri ini.
  \item Bab~24 (panas internal Bumi): pembangkit listrik panas bumi Nesjavellir, Islandia, foto karya Gretar Ívarsson, disunting oleh Fir0002, domain publik (Wikimedia Commons).
  \item Bab~24: aliran panas dari Davies dan Davies (2010); panas radiogenik dari pengukuran geoneutrino KamLAND (2011) dan pelepasan panas keempat isotop pada masa kini; produksi panas batuan dari Vilà dan Fernández (2010).
  \item Bab~25 (penanggalan radiometrik dan umur Bumi): William Thomson, Lord Kelvin, 1906, foto karya Russell and Sons, London, Smithsonian Institution Libraries, tanpa pembatasan hak cipta yang diketahui; meteorit besi Canyon Diablo, foto karya Geoffrey Notkin, CC~BY-SA~2.5 (keduanya Wikimedia Commons).
  \item Bab~25: isotop timbal meteorit dari Patterson (1956); konstanta peluruhan rubidium-87 yang direkomendasikan IUGS dan IUPAC (2022); sejarah perkiraan-perkiraan dari Holmes (1913, 1927), Patterson (1956), dan U.S. Geological Survey; Perbukitan Varden, wilayah rekaan seri ini.
  \item Bab~26 (Bumi purba dan atmosfernya): formasi besi berpita, Dales Gorge, foto karya Graeme Churchard, CC~BY~2.0; stromatolit di Teluk Shark, foto karya Paul Harrison, CC~BY-SA~3.0; lapisan merah dekat Sedona, foto karya Wendy from Pennsylvania, CC~BY-SA~2.0 (semuanya Wikimedia Commons).
  \item Bab~26: luminositas Matahari muda dari Gough (1981); batas dan warna eon dari International Chronostratigraphic Chart.
  \item Bab~27 (membaca iklim masa lalu): gelembung udara dalam es Antarktika, foto CSIRO Atmospheric Research, CC~BY~3.0; foraminifera planktonik, foto karya Marci Robinson, USGS, domain publik; butir-butir polen, Dartmouth Electron Microscope Facility, domain publik (semuanya Wikimedia Commons).
  \item Bab~27: suhu EPICA Dome~C (Jouzel dan kawan-kawan, 2007) dan karbon dioksida inti es Antarktika (Bereiter dan kawan-kawan, 2015) melalui NOAA Paleoclimatology; tumpukan LR04 (Lisiecki dan Raymo, 2005); relief ETOPO1 (NOAA); NASA GISS (GISTEMP v4); periode acuan dari Laporan Penilaian Keenam IPCC, Kelompok Kerja~I, bab~2.
  \item Bab~28 (perubahan iklim): Observatorium Mauna Loa; foto karya Brian Vasel, NOAA, domain publik (Wikimedia Commons).
  \item Bab~28: Charles David Keeling dan John Chin di Mauna Loa, 1997; foto karya John Miller, NOAA, domain publik (Wikimedia Commons).
  \item Bab~28: data dari NOAA Global Monitoring Laboratory dan Scripps Institution of Oceanography (karbon dioksida), NASA GISS (GISTEMP v4), NOAA Laboratory for Satellite Altimetry (permukaan laut), Bereiter dan kawan-kawan (2015) melalui NOAA Paleoclimatology (inti es), dan Laporan Penilaian Keenam IPCC, Kelompok Kerja~I (gaya radiatif, atribusi, proyeksi).
\end{itemize}
